\scanpage{13}
! On obtient donc des bijections inverses l'une de l'autre.
\diagram{\begin{tikzcd}[ampersand replacement=\&]
\pi_0(U\otimes\Lambda)\arrow[r,shift left,"\exp"] \& \gamma_0(U\otimes_k\Lambda)\arrow[l,shift left,"\log"]
\end{tikzcd}}
où les indices $0$ indiquent qu'on
prend les primitifs resp$^{\mathrm t}$ multiplicatifs
``relatifs'' à l'augmentation $U\otimes\Lambda\longrightarrow U$.
[\unsure{passage|passant} à la limite projective \inserted{sur $\Lambda$}, on déduit
un résultat analogue si $\Lambda$ est une
limite projective de $k$-algèbres infinitésimales]

Notons que \deleted{si $\Lambda$ est plat, \unsure{ou} si} \inserted{l'on a d'ailleurs} un
\unsure{homom|homom.} naturel de $\Lambda$-modules
\[
\pi(U)\otimes_k I(\Lambda)\xrightarrow{\alpha}\pi_0(U\otimes_k\Lambda)
\]
qui est un isom si $\Lambda$ est $k$-plat,
ou si
\[
0\longrightarrow\deleted{\unsure{\gamma}}\pi(U)\longrightarrow\deleted{\unsure{I(U)}}\unsure{U|I(U)}\xrightarrow{\delta} I(U)\otimes I(U)
\]
est une suite \uuline{\unsure{purement|pûrement} exacte}. On
en déduit donc un hom canonique
\[
\pi(U)\otimes_k I(\Lambda)\xrightarrow{\exp\circ\alpha}1+\pi_0(U\otimes_k\Lambda)
\]
qui est un isom. dans les cas
qu'on vient de préciser.

Conséquences : Supposons $k$ corps, $U$ une hyperalgèbre
\deleted{\unsure{au sens de}} \uline{de Cartier} \unsure{$\rho$}
Alors $U$ (avec $\Delta$ \uline{cocomm.} et \uline{ass.}) est connu
quand on connait le foncteur $\pi_0(U\otimes_k\Lambda)$. $\rho$

% margin note written sideways
\margin{$\rho$ [i.e. $\Delta$ est cocomm. et assoc, \deleted{et \unsure{il y a}} $I(U)=\bigcup U_i$, avec $\delta U_i\subset U\otimes U_i$ et $U_i$ de dim. finie, enfin $\gamma(U)=\{e\}$]}
% margin note written sideways
\margin{$\rho$ \unsure{En dernière analyse|La donnée d'entre} \unsure{tout} est connu quand on connait $\inserted{\mathfrak g=}\pi(U)\deleted{\unsure{\otimes}}$ (formule de Hausdorff-Campbell), donc comme par Poincaré--Witt on a $\mathfrak g\simeq\pi(U(\mathfrak g))$ on voit que $\unsure{U_h|U}\simeq U(\mathfrak g)$ où $\mathfrak g=\pi(U)$ : \uline{th. de} \uline{Cartier} \inserted{En} particulier, si $\mathfrak g$ est de dim. finie, le dual de $\unsure{\mathfrak g|U}$ est \uline{un ann. \unsure{local}} \uline{régulier}. \uline{Cor} Un groupe alg. en car\unsure{q.} $0$ est \unsure{simple|lisse} sur $k$.}
% archive sheet number: 12

